% Requires: amsmath, amssymb, booktabs, float, algorithm, algpseudocode, bm

\section{Centralized Multi-Period Optimal Power Flow}
\label{sec:copf}

We study the day-ahead operation of an unbalanced radial distribution
network (DN) hosting distributed energy resources (DERs) and battery
energy storage systems (BESS). Over a scheduling horizon
$\mathcal{T}=\{1,\dots,T\}$, the operator minimizes the cost of energy
imported at the substation subject to the network physics and all
device limits. Because each battery couples consecutive timesteps
through its state of charge (SOC), the resulting program is a single,
temporally coupled multi-period OPF. Section~\ref{ssec:notation} fixes
the notation; Section~\ref{ssec:nonconvex} states the exact non-convex
model; Section~\ref{ssec:mpisocp} develops the iterative SOCP scheme
that solves it; and Section~\ref{ssec:linear} gives the linear
counterpart. These models are the building blocks placed inside each
temporal window of the decomposition in Section~\ref{sec:dddp}.

\subsection{Notation and Network Model}
\label{ssec:notation}

The DN is a directed graph $G=(\mathcal{N},\mathcal{E})$ with node set
$\mathcal{N}$ and branch set $\mathcal{E}$; $\Phi_j$ and $\Phi_{ij}$
collect the phases present at bus $j$ and branch $(i,j)$, superscript
$p$ indexes the phase, and $t$ indexes time. For branch $(i,j)$ and
phase $p$ at time $t$, $P^{pp}_{ij,t}$ and $Q^{pp}_{ij,t}$ are the
active and reactive power flows. At node $j$, $p^{p}_{L,j,t}$ and
$q^{p}_{L,j,t}$ are the load demands and $p^{p}_{D,j,t}$, $q^{p}_{D,j,t}$
the DER injections, bounded by the inverter rating
$S^{p}_{\mathrm{DR},j}$. The BESS at node $j$ charges and discharges
with powers $P^{p}_{c,j,t}$, $P^{p}_{d,j,t}$, stores energy (SOC)
$B^{p}_{j,t}$, has rated power $P_{\mathrm{BR},j}$, efficiencies
$\eta_c,\eta_d$, and SOC bounds
$[\mathrm{soc}_{\min},\mathrm{soc}_{\max}]$. We use the squared voltage
magnitude $\upsilon^{p}_{j,t}=|V^{p}_{j,t}|^2$, the squared current
magnitude $l^{pp}_{ij,t}=|I^{pp}_{ij,t}|^2$, the cross-phase current
product $l^{pq}_{ij,t}=|I^{pp}_{ij,t}|\,|I^{qq}_{ij,t}|$, and the
current-angle difference $\delta^{pq}_{ij,t}$. Line impedance and
resistance are $z^{pq}_{ij}$, $R^{pq}_{ij}$, and voltages are confined
to $[V_{\min},V_{\max}]$. Table~\ref{tab:notation} summarizes the
principal symbols.

\begin{table}[H]
  \centering
  \caption{Key Notation}
  \label{tab:notation}
  \begin{tabular}{ll}
    \toprule
    Symbol & Description \\
    \midrule
    \(G=(\mathcal{N},\mathcal{E})\) & Directed graph of nodes and branches \\
    \(\Phi_j,\,\Phi_{ij}\) & Set of phases at bus \(j\) and branch \((i,j)\) \\
    \(P^{pp}_{ij,t},\,Q^{pp}_{ij,t}\) & Active/reactive flow on \((i,j)\), phase \(p\), time \(t\) \\
    \(p^{p}_{L,j,t},\,q^{p}_{L,j,t}\) & Active/reactive demand at node \(j\), phase \(p\) \\
    \(p^{p}_{D,j,t},\,q^{p}_{D,j,t}\) & DER active/reactive injection at node \(j\) \\
    \(S^{p}_{\mathrm{DR},j}\) & DER inverter apparent-power rating \\
    \(P^{p}_{c,j,t},\,P^{p}_{d,j,t}\) & BESS charge/discharge active power \\
    \(B^{p}_{j,t}\) & BESS state of charge (SOC) \\
    \(P_{\mathrm{BR},j}\) & BESS rated active-power capacity \\
    \(\eta_c,\,\eta_d\) & BESS charge/discharge efficiencies \\
    \(\mathrm{soc}_{\min},\,\mathrm{soc}_{\max}\) & BESS SOC bounds \\
    \(V^{p}_{j,t}\) & Complex voltage \((|V^{p}_{j,t}|\angle\theta^{p}_{j,t})\) \\[2pt]
    \(\upsilon^{p}_{j,t}\) & Squared voltage magnitude \((\upsilon^{p}_{j,t}=|V^{p}_{j,t}|^2)\) \\[2pt]
    \(I^{pp}_{ij,t}\) & Complex current \((|I^{pp}_{ij,t}|\angle\delta^{pp}_{ij,t})\) \\[2pt]
    \(\delta^{pq}_{ij,t}\) & Current-angle difference \((\delta^{pp}_{ij,t}-\delta^{qq}_{ij,t})\) \\[2pt]
    \(l^{pp}_{ij,t}\) & Squared current magnitude \((|I^{pp}_{ij,t}|^2)\) \\[2pt]
    \(l^{pq}_{ij,t}\) & Current-magnitude product \((|I^{pp}_{ij,t}|\,|I^{qq}_{ij,t}|)\) \\[2pt]
    \(z^{pq}_{ij},\,R^{pq}_{ij}\) & Impedance/resistance of line \((i,j)\), \(pq\in\Phi_{ij}\) \\
    \bottomrule
  \end{tabular}
\end{table}

\subsection{Non-Convex Multi-Period Formulation}
\label{ssec:nonconvex}

The centralized multi-period three-phase OPF is stated compactly by
\eqref{eq0}--\eqref{lintermsoc}, and its terms are described immediately
afterward. The problem reads
\begin{equation}\label{eq0}
\small
    \min_{\mathcal{Y}}\; F(\mathcal{Y}),
\end{equation}
\begin{equation}\label{eq1}
\begin{small}
\begin{aligned}
F(\mathcal{Y})=\sum_{t=1}^{T} \Bigg[
& C_{t}\,P_{subs,t}
+ \alpha \sum_{j \in \mathcal{B}} \left\{ (1 - \eta_c) P_{C_j,t}
+ \left( \tfrac{1}{\eta_d} - 1 \right) P_{D_j,t} \right\}
\Bigg],
\end{aligned}
\end{small}
\end{equation}
subject to, for every branch $(i,j)\in\mathcal{E}$, phase
$p\in\Phi_{ij}$, and time $t\in\mathcal{T}$,
\begin{equation}\label{eq2}
\begin{split}
P_{ij,t}^{pp}
= \sum_{(j,k)\in E} P_{jk,t}^{pp} + p_{L_j,t}^p
- p_{D_j,t}^p - P_{d_j,t}^p + P_{c_j,t}^p \\ + \sum_{q \in \Phi_j} l_{ij,t}^{pq}
\left( r_{ij}^{pq} \cos(\delta_{ij,t}^{pq}) + x_{ij}^{pq} \sin(\delta_{ij,t}^{pq}) \right),
\end{split}
\end{equation}
\begin{equation}\label{eq3}
\begin{split}
Q_{ij,t}^{pp}
= & \sum_{(j,k)\in E} Q_{jk,t}^{pp} + q_{L_j,t}^p - q_{D_j,t}^p \\
& + \sum_{q \in \Phi_j} l_{ij,t}^{pq}
\left( x_{ij}^{pq} \cos(\delta_{ij,t}^{pq}) - r_{ij}^{pq} \sin(\delta_{ij,t}^{pq}) \right),
\end{split}
\end{equation}
\begin{equation}\label{eq4}
\begin{split}
v_{j,t}^p =  & v_{i,t}^p
- \sum_{q \in \Phi_j} 2\,\mathbb{R}\!\left[ S_{ij,t}^{pq} \left( z_{ij}^{pq} \right)^* \right]- \\ & \sum_{\substack{q_1, q_2 \in \Phi_j}}
2\,\mathbb{R}\!\left[ z_{ij}^{pq_1} l_{ij,t}^{q_1 q_2}
\left( \angle(\delta_{ij,t}^{q_1 q_2}) \right)
\left( z_{ij}^{pq_2} \right)^* \right],
\end{split}
\end{equation}
\begin{equation}\label{eq5}
\left( P_{ij,t}^{pp} \right)^2 + \left( Q_{ij,t}^{pp} \right)^2
= v_{i,t}^p \, l_{ij,t}^{pp},
\end{equation}
\begin{equation}\label{eq6}
\left( l_{ij,t}^{pq} \right)^2 = l_{ij,t}^{pp} \, l_{ij,t}^{qq},
\end{equation}
\begin{equation}\label{linqder}
{ - \sqrt {{{\left( {S_{D\,R_j}^p} \right)}^2} - {{\left( {p_{D_j,t}^p} \right)}^2}} \le q_{D_j,t}^p \le \sqrt {{{\left( {S_{D\,R_j}^p} \right)}^2} - {{\left( {p_{D_j,t}^p} \right)}^2}}},
\end{equation}
\begin{equation}\label{linv}
{\left( {{V_{min}}} \right)^2} \le {\upsilon}_j^p \le {\left( {{V_{max}}} \right)^2},
\end{equation}
\begin{equation}\label{linsocevol}
B_{j,t}^{p} = B_{j,t-1}^{p} + \Delta t\, \eta_c P_{c_j,t}^{p} - \Delta t\, (1/\eta_d) P_{d_j,t}^{p},
\end{equation}
\begin{equation}\label{linPch}
P_{cj,t}^{p} \le P_{BR_j},
\end{equation}
\begin{equation}\label{linPdis}
P_{dj,t}^{p} \le P_{BR_j},
\end{equation}
\begin{equation} \label{linsoclim}
soc_{min}B_{R,j} \le B_{j,t}^p \le  soc_{max}B_{R,j},
\end{equation}
\begin{equation}\label{lintermsoc}
B_{j,0}^{p} = B_{j,T}^{p}.
\end{equation}

The decision set $\mathcal{Y}$ collects all line flows, voltages, DER
injections, and BESS variables over the horizon. In the
objective~\eqref{eq1}, the first term prices substation imports ($C_t$
in $\$/\mathrm{kWh}$, $P_{subs,t}$ in kW) and the second penalizes BESS
conversion losses with weight $\alpha>0$; as shown in~\cite{Nazir}, this
penalty precludes simultaneous charging and discharging without
introducing binary variables, keeping the model continuous.
Constraints~\eqref{eq2}--\eqref{eq3} enforce three-phase active and
reactive nodal balance, the trailing summations accounting for series
line losses; \eqref{eq4} is the three-phase Kirchhoff voltage law across
branch $(i,j)$; and \eqref{eq5}--\eqref{eq6} are the exact
power--voltage--current relation and the cross-phase current-coupling
identity of the branch-flow model. The DER reactive capability is
box-bounded by~\eqref{linqder} and bus voltages by~\eqref{linv}. The
BESS is governed by the SOC recursion~\eqref{linsocevol}, the
charge/discharge power caps~\eqref{linPch}--\eqref{linPdis}, the energy
bounds~\eqref{linsoclim}, and the cyclic terminal
condition~\eqref{lintermsoc}. Crucially, \eqref{linsocevol} is the sole
inter-temporal coupling: every $B^{p}_{j,t}$ with $t\ge1$ is a decision
variable, while at the first step $t=1$ the preceding state is fixed to
the given initial SOC, $B^{p}_{j,0}$, which anchors the entire horizon,
and \eqref{lintermsoc} restores that level at $t=T$.

The formulation is non-convex on two counts. Treating the current
angles $\delta^{pq}_{ij,t}$ as variables makes
\eqref{eq2}--\eqref{eq4} strongly nonlinear; however, in distribution
networks these angles deviate only marginally from their base-case
values~\cite{XX}, so we fix $\delta^{pq}_{ij,t}$ to the angles obtained
from a base-case OpenDSS power flow, held constant over all timesteps.
This renders \eqref{eq2}--\eqref{eq4} linear. The only residual
non-convexities are then the equalities~\eqref{eq5}--\eqref{eq6}, whose
bilinear right-hand sides are addressed next.

\subsection{MPISOCP: Iterative SOCP Convexification}
\label{ssec:mpisocp}

We convexify the remaining non-convex equalities by a Multi-Period
Iterative Second-Order Cone Programming (MPISOCP) scheme, whose
ingredients are collected in \eqref{eqpvc}--\eqref{eqdircmc} and
explained thereafter. The equalities~\eqref{eq5}--\eqref{eq6} are first
relaxed into second-order cones,
\begin{equation}\label{eqpvc}
\left( P_{ij,t}^{pp} \right)^2 + \left( Q_{ij,t}^{pp} \right)^2
\leq v_{i,t}^p \, l_{ij,t}^{pp},
\end{equation}
\begin{equation}\label{eqcmc}
\left( l_{ij,t}^{pq} \right)^2 \leq l_{ij,t}^{pp} \, l_{ij,t}^{qq},
\end{equation}
their (non-positive) relaxation gaps are quantified by the residuals
\begin{equation}\label{erroreqpvc}
\left( E1_{ij,t}^{pp} \right) = \left( P_{ij,t}^{pp} \right)^2 + \left( Q_{ij,t}^{pp} \right)^2
- v_{i,t}^p \, l_{ij,t}^{pp},
\end{equation}
\begin{equation}\label{erroreqcmc}
\left( E2_{ij,t}^{pq} \right) = \left( l_{ij,t}^{pq} \right)^2 - l_{ij,t}^{pp} \, l_{ij,t}^{qq},
\end{equation}
these residuals are driven toward zero across iterations by
\begin{equation}\label{erroreqpvcred}
(E1_{ij,t}^{pp})^{m+1} \geq \tau\, (E1_{ij,t}^{pp})^{m},
\end{equation}
\begin{equation}\label{erroreqcmcred}
(E2_{ij,t}^{pq})^{m+1} \geq \tau\, (E2_{ij,t}^{pq})^{m},
\end{equation}
and, after a first-order Taylor expansion about the iterate-$m$
operating point, the tractable directional cuts are
\begin{equation}\label{eqdirpvc}
\begin{split}
2P_{ij,t}^{pp(m)} P_{ij,t}^{pp(m+1)}
&+ 2Q_{ij,t}^{pp(m)} Q_{ij,t}^{pp(m+1)}
- l_{ij,t}^{pp(m)} v_{i,t}^{p(m+1)} \\&
- v_{i,t}^{p(m)} l_{ij,t}^{pp(m+1)}
\geq (\tau + 1)\, E1_{ij,t}^{pp(m)},
\end{split}
\end{equation}
\begin{equation}\label{eqdircmc}
\begin{split}
2l_{ij,t}^{pq(m)}  l_{ij,t}^{pq(m+1)}
- l_{ij,t}^{pp(m)}  l_{ij,t}^{qq(m+1)}
- l_{ij,t}^{qq(m)}  l_{ij,t}^{pp(m+1)} \\
\geq (\tau + 1)\, E2_{ij,t}^{pq(m)}.
\end{split}
\end{equation}

The inequalities~\eqref{eqpvc}--\eqref{eqcmc} enlarge the feasible set
and make each stage a convex conic program; the relaxation is exact only
when both cones are tight, i.e.\ when the residuals
\eqref{erroreqpvc}--\eqref{erroreqcmc} vanish. Since these residuals are
non-positive by construction, \eqref{erroreqpvcred}--\eqref{erroreqcmcred}
require each one at iteration $m{+}1$ to reach at least a fraction
$\tau\in[0,1]$ of its previous value, pushing it toward zero; the
right-hand sides are constants known after iteration $m$, but the
left-hand sides remain bilinear. Linearizing them yields the linear
directional cuts~\eqref{eqdirpvc}--\eqref{eqdircmc}, which steer the
conic solution onto the exact manifold and shrink the gap monotonically.
Algorithm~\ref{alg:isocp} iterates this process: it solves the SOCP
relaxation with the current cuts, re-linearizes at the new iterate, and
terminates once the maximum residual falls below $\varepsilon$ or the
iteration cap $M_{\max}$ is reached.

\begin{algorithm}[t]
\caption{Multi-Period Iterative SOCP (MPISOCP).}
\label{alg:isocp}
\begin{algorithmic}[1]
\Require $\tau\!\in\![0,1]$, tolerance $\varepsilon$, cap $M_{\max}$.
\State Solve the SOCP relaxation
\eqref{eq0}--\eqref{eq4},\,\eqref{eqpvc},\,\eqref{eqcmc}, \eqref{linqder}--\eqref{lintermsoc}; obtain
$\mathcal{Y}^{(0)}\!=\!(P^{(0)},Q^{(0)},v^{(0)},l^{(0)})$.
\State Compute $E1^{pp(0)}_{ij,t}$, $E2^{pq(0)}_{ij,t}$
from~\eqref{erroreqpvc}--\eqref{erroreqcmc} and
$E^{(0)}\!\leftarrow\!\max\{|E1^{pp(0)}|,|E2^{pq(0)}|\}$.
\For{$m=0,1,\dots,M_{\max}\!-\!1$}
    \If{$E^{(m)}\!\le\!\varepsilon$} \Return $\mathcal{Y}^{(m)}$. \EndIf
        \State Linearize the residuals at $\mathcal{Y}^{(m)}$ to form the
directional cuts~\eqref{eqdirpvc}--\eqref{eqdircmc}.
\State Solve \eqref{eq0}--\eqref{eq4},\,\eqref{eqpvc},\,\eqref{eqcmc},\eqref{linqder}--\eqref{lintermsoc}, augmented with the cuts~\eqref{eqdirpvc}--\eqref{eqdircmc}; obtain
$\mathcal{Y}^{(m+1)}$ and update $E^{(m+1)}$.
\EndFor
\State \Return $\mathcal{Y}^{(M_{\max})}$.
\end{algorithmic}
\end{algorithm}

\subsection{Linear Multi-Period Formulation}
\label{ssec:linear}

When speed is favored over loss accuracy, the quadratic loss terms are
dropped, yielding a fully linear multi-period OPF. It retains the
objective~\eqref{eq1} and the DER/BESS
constraints~\eqref{linqder}--\eqref{lintermsoc}, and replaces
\eqref{eq2}--\eqref{eq4} by their loss-free counterparts
\begin{equation}\label{linPij}
P_{ij,t}^{pp} = \sum_{(j,k)\in E} {P_{jk,t}^{pp}} + p_{{L_j,t}}^p - p_{D_j,t}^p - P_{d_j,t}^p + P_{c_j,t}^p,
\end{equation}
\begin{equation}\label{linQij}
Q_{ij,t}^{pp} = \sum_{(j,k)\in E} {Q_{jk,t}^{pp}} + q_{{L_j,t}}^p - q_{Dj,t}^p,
\end{equation}
\begin{equation}\label{linkvl}
{\upsilon }_{j,t}^p = {\upsilon}_{i,t}^p - \sum\limits_{q \in {\phi _j}} {2\left[ {S_{ij,t}^{pq}\left( {z_{ij}^{pq}} \right)^*} \right]}.
\end{equation}
Equations~\eqref{linPij}--\eqref{linQij} are the lossless nodal balances
and \eqref{linkvl} the lossless voltage drop. This model has no bilinear
terms, needs no convexification, and is solved directly as a single
linear program; it is the fast alternative to MPISOCP inside the
temporal decomposition of Section~\ref{sec:dddp}.

\section{Deterministic Dual Dynamic Programming (DDDP)}
\label{sec:dddp}

The centralized problem of Section~\ref{sec:copf} solves the whole
horizon $\mathcal{T}=\{1,\dots,T\}$ as one program, whose timesteps are
tied together solely by the SOC recursion~\eqref{linsocevol}. Even this
single coupling makes the monolithic solve time grow super-linearly with
$T$, which is prohibitive for long-horizon, feeder-scale studies. We
therefore decompose the horizon in time: the schedule is split into a
chain of shorter window sub-problems that communicate only across their
boundaries, and the centralized optimum is recovered by a parallel
chained Benders scheme we call Deterministic Dual Dynamic Programming
(DDDP).

The windows couple through one state variable---the boundary SOC---but
DDDP exchanges information about it in two passes. The \emph{forward}
pass carries the SOC value: each window hands its terminal SOC to the
next as an initial condition. The \emph{backward} pass carries a Benders
cut: each window returns to its predecessor a linear under-estimate of
the downstream cost incurred by that SOC. Because both passes read only
quantities from the \emph{previous} iteration, all $P$ windows are solved
simultaneously.

\subsection{Window Construction and Boundary Variables}
\label{ssec:windows}

We partition $\mathcal{T}$ into $P$ \emph{disjoint} contiguous windows
$\{\mathcal{W}_i\}_{i=1}^{P}$ of length $b$, where window $i$ spans
$\mathcal{W}_i=\{w^{s}_i,\dots,w^{e}_i\}$ with
\begin{equation}
\label{eq:windows}
w^{s}_i = (i\!-\!1)b+1, \qquad w^{e}_i = ib,
\qquad w^{s}_1 = 1,\; w^{e}_P = T .
\end{equation}
The windows carry no look-ahead overlap; the influence of future stages
is captured exactly through Benders cuts, which later makes the stitched
primal cost exact (Section~\ref{ssec:bounds}). The only primal variable
coupling adjacent windows is the terminal SOC of window $i$, defined over
the battery buses $\mathcal{B}$ as
\begin{equation}
\label{eq:state_definition}
\bm{x}_i \;\triangleq\; \big(B^{p}_{j,\,w^{e}_i}\big)_{j\in\mathcal{B},\,p\in\Phi_j},
\qquad
\bm{x}_0 \;\triangleq\; \big(B^{p}_{j,0}\big)_{j\in\mathcal{B},\,p\in\Phi_j}.
\end{equation}
The forward pass propagates the \emph{value} of $\bm{x}_i$, handed to
window $i{+}1$ as its fixed initial SOC; the backward pass propagates
\emph{dual information} about the same $\bm{x}_i$, delivered as the
Benders cut of Section~\ref{ssec:cuts}. For the first window the initial
SOC is $\bm{x}_0$, the given value that anchors the centralized horizon
and never changes.

\subsection{Window Sub-Problem}
\label{ssec:subproblem}

Each window $i$ solves the chosen centralized OPF---the MPISOCP model of
Section~\ref{ssec:mpisocp} or the linear model of
Section~\ref{ssec:linear}---over the reduced horizon $\mathcal{W}_i$,
augmented by a scalar epigraph variable $\theta_i\ge0$ that bounds the
optimal cost of the downstream windows $i{+}1,\dots,P$. Let
$\mathcal{Y}_i$ be the local decision vector, $F(\mathcal{Y}_i)$ the
centralized objective~\eqref{eq1} summed over $t\in\mathcal{W}_i$, and
$\mathcal{F}_i(\hat{\bm{x}}_{i-1})$ the window feasible set parameterized
by the incoming SOC. That incoming SOC is a fixed \emph{parameter},
distinguished from the decision variable $\bm{x}_{i-1}$ by a hat; at
outer iteration $k$ it takes the trial value
\begin{equation}
\label{eq:trial}
\hat{\bm{x}}_{i-1}^{(k)} \;\triangleq\; \bm{x}_{i-1}^{(k-1)},
\end{equation}
the terminal SOC produced by window $i{-}1$ at the previous iteration
(with $\hat{\bm{x}}_0^{(k)}\equiv\bm{x}_0$ for all $k$). The window-$i$
sub-problem at iteration $k$, denoted $[\mathrm{SP}_i^{(k)}]$, is
\begin{subequations}
\label{eq:subproblem}
\begin{align}
V_i^{(k)}\big(\hat{\bm{x}}_{i-1}^{(k)}\big) \;\triangleq\;
\min_{\mathcal{Y}_i,\,\theta_i}\;\;
& F(\mathcal{Y}_i)\;+\;\theta_i
\label{eq:sub_obj}\\[2pt]
\text{s.t.}\quad
& \mathcal{Y}_i \,\in\, \mathcal{F}_i(\hat{\bm{x}}_{i-1}^{(k)}),
\label{eq:sub_feas}\\[2pt]
& \theta_i \;\ge\; \alpha_i^{(\ell)} \;+\; \bm{\beta}_i^{(\ell)\top}\, \bm{x}_i,
   \quad \ell = 1,\dots,k{-}1,
\label{eq:sub_cuts}\\[2pt]
& \theta_i \;\ge\; 0 .
\label{eq:sub_thpos}
\end{align}
\end{subequations}
The augmented objective~\eqref{eq:sub_obj} is the window value function:
the local operating cost plus the value-to-go, evaluated at the trial
point $\hat{\bm{x}}_{i-1}^{(k)}$. Constraint~\eqref{eq:sub_feas} is the
complete MPISOCP or linear model over $\mathcal{W}_i$, so the SOC
recursion~\eqref{linsocevol} holds for every $t\in\mathcal{W}_i$;
and~\eqref{eq:sub_cuts} is the pool of $k{-}1$ optimality cuts received
from window $i{+}1$, written in the terminal-SOC \emph{variable}
$\bm{x}_i$, with $\theta_P\equiv0$ for the last window.

The trial SOC enters $[\mathrm{SP}_i^{(k)}]$ only through the initial
condition of the recursion~\eqref{linsocevol}: window $i$ fixes its
pre-window state at $t=w^{s}_i$ to $\hat{\bm{x}}_{i-1}^{(k)}$, exactly as
the centralized problem fixes its pre-horizon state to $\bm{x}_0$ at
$t=1$. Consequently the dual multiplier of this fixed initial condition
equals the gradient of the value function at the trial point,
\begin{equation}
\label{eq:dual_sens}
\bm{\beta}_i^{(k)} \;=\;
\frac{\partial V_i^{(k)}}{\partial \bm{x}_{i-1}}\bigg|_{\hat{\bm{x}}_{i-1}^{(k)}} ,
\end{equation}
the sensitivity used to build the backward Benders cut of
Section~\ref{ssec:cuts}.

\subsection{Cut Generation: A Walk-Through of Window~$i$}
\label{ssec:cuts}

To make the two-pass exchange concrete, we trace one outer iteration $k$
for an interior window $i\in\{2,\dots,P{-}1\}$; the boundary windows
differ only in that window $1$ uses the fixed trial
$\hat{\bm{x}}_0^{(k)}\equiv\bm{x}_0$ and window $P$ has
$\theta_P\equiv0$. Window $i$ executes five steps.
\begin{enumerate}\itemsep2pt
\item[\textbf{S1.}] \emph{Receive forward state.} Fix the incoming SOC
to the trial value $\hat{\bm{x}}_{i-1}^{(k)}$ of~\eqref{eq:trial}, the
terminal SOC delivered by window $i{-}1$ at iteration $k{-}1$.
\item[\textbf{S2.}] \emph{Receive backward cut.} Append the optimality
cut generated by window $i{+}1$ at iteration $k{-}1$ to the
pool~\eqref{eq:sub_cuts}.
\item[\textbf{S3.}] \emph{Solve.} Solve $[\mathrm{SP}_i^{(k)}]$ for the
value $V_i^{(k)}$, the primal $\mathcal{Y}_i^{(k)}$ (which contains the
new terminal SOC $\bm{x}_i^{(k)}$), and the boundary dual
$\bm{\beta}_i^{(k)}$ of~\eqref{eq:dual_sens}.
\item[\textbf{S4.}] \emph{Emit forward state.} Pass $\bm{x}_i^{(k)}$ to
window $i{+}1$ as its next trial $\hat{\bm{x}}_i^{(k+1)}$.
\item[\textbf{S5.}] \emph{Emit backward cut.} Return to window $i{-}1$ an
optimality cut on its value-to-go variable $\theta_{i-1}$, written in
that window's terminal-SOC variable $\bm{x}_{i-1}$,
\begin{equation}
\label{eq:cut_build}
\theta_{i-1} \;\ge\;
\underbrace{V_i^{(k)} - \bm{\beta}_i^{(k)\top}\hat{\bm{x}}_{i-1}^{(k)}}_{\alpha_{i-1}^{(k)}}
\;+\;
\underbrace{\bm{\beta}_i^{(k)\top}}_{\bm{\beta}_{i-1}^{(k)\top}}\,
\bm{x}_{i-1}.
\end{equation}
\end{enumerate}
The cut~\eqref{eq:cut_build} is the supporting hyperplane of the
downstream value function $V_i$ at the trial point: its slope is the
sensitivity $\bm{\beta}_i^{(k)}$ of~\eqref{eq:dual_sens} and its
intercept anchors it to $(\hat{\bm{x}}_{i-1}^{(k)},V_i^{(k)})$. Since
$V_i$ is convex in its boundary argument under the SOCP relaxation, the
cut is a valid under-estimator of $\theta_{i-1}$, and the accumulated
cuts form a piecewise-linear surrogate that tightens monotonically.
Because every window reads only trial states and cuts from the previous
iteration (S1--S2), the $P$ sub-problems are decoupled within an
iteration and solved in parallel, exchanging just one forward SOC vector
$\bm{x}_i$ and one cut $(\alpha,\bm{\beta})$ per window per iteration. The
complete procedure is summarized in Algorithm~\ref{alg:dddp}.

\subsection{Global Bounds and Convergence}
\label{ssec:bounds}

Because the windows are disjoint and exhaustive, each outer iteration
$k$ exposes a global lower and upper bound at negligible cost,
\begin{align}
\underline{Z}^{(k)} \;&\triangleq\; V_1^{(k)}
\;=\; F(\mathcal{Y}_1^{(k)}) + \theta_1^{(k)},
\label{eq:LB}\\[2pt]
\overline{Z}^{(k)} \;&\triangleq\; \sum_{i=1}^{P} F(\mathcal{Y}_i^{(k)}),
\label{eq:UB}
\end{align}
and convergence is declared when either of the two criteria
\begin{align}
\textbf{(C1)}\;\;\;
\frac{\overline{Z}^{(k)} - \underline{Z}^{(k)}}
     {\max(1,|\overline{Z}^{(k)}|)} &\;<\; \varepsilon_{\text{gap}},
\label{eq:cv_gap}\\[2pt]
\textbf{(C2)}\;\;\;
\max\!\Big\{\!
\tfrac{|\underline{Z}^{(k)}{-}\underline{Z}^{(k-1)}|}{\max(1,|\underline{Z}^{(k)}|)},\;
\tfrac{|\underline{Z}^{(k-1)}{-}\underline{Z}^{(k-2)}|}{\max(1,|\underline{Z}^{(k)}|)}\!
\Big\} &\;<\; \varepsilon_{\text{lb}}
\label{eq:cv_lb}
\end{align}
is met. In~\eqref{eq:LB}, $\underline{Z}^{(k)}$ is a global lower bound
on the centralized optimum $Z^\star$: $\theta_1$ under-estimates the
true value-to-go of windows $2,\dots,P$, so $V_1^{(k)}\le Z^\star$ for
all $k$ and is non-decreasing under standard Benders
monotonicity~\cite{XX}. In~\eqref{eq:UB}, $\overline{Z}^{(k)}$ is the
\emph{exact} primal cost of the stitched trajectory: with zero overlap
each timestep is optimized by exactly one window, so the stage costs add
up with no double-counting. Of the two rules, (C1) is the classical
Benders duality gap, while (C2) is a two-step lower-bound stability test
that guards against the mild primal degeneracy at window
boundaries---which makes $\overline{Z}^{(k)}$ oscillate by a few cuts
even after the lower bound has plateaued---and empirically triggers
first on feeders with many co-located batteries.

\begin{algorithm}[t]
\caption{Parallel DDDP.}
\label{alg:dddp}
\begin{algorithmic}[1]
\Require horizon $T$, partitions $P$, initial SOC $\bm{x}_0$,
         tolerances $\varepsilon_{\text{gap}},\varepsilon_{\text{lb}}$,
         max iterations $K$.
\State Build disjoint windows $\{\mathcal{W}_i\}_{i=1}^{P}$
       via~\eqref{eq:windows}.
\State Spawn $P$ persistent workers; each builds its local model
       $[\mathrm{SP}_i^{(0)}]$ once with an empty cut-pool.
\State Initialize $\hat{\bm{x}}_i^{(1)}\leftarrow\bm{x}_0$ and cut inbox
       $\mathcal{C}_i\leftarrow\emptyset$ for all $i$.
\For{$k = 1, 2, \dots, K$}
   \State \textbf{(Parallel solve)}\;
          \textbf{parfor} $i = 1, \dots, P$ \textbf{do}
   \State \quad drain inbox: append cuts $\mathcal{C}_i$ to the pool of
          $[\mathrm{SP}_i^{(k)}]$, then set $\mathcal{C}_i\leftarrow\emptyset$
   \State \quad solve $[\mathrm{SP}_i^{(k)}]$ at $\hat{\bm{x}}_{i-1}^{(k)}$,
          obtaining $V_i^{(k)},\bm{x}_i^{(k)},\bm{\beta}_i^{(k)}$
   \State \textbf{end parfor}
   \State \textbf{(Cut assembly)} \For{$i = 2, \dots, P$}
       \State $\alpha \leftarrow V_i^{(k)} -
              \bm{\beta}_i^{(k)\top}\hat{\bm{x}}_{i-1}^{(k)}$;\;\;
              $\bm{\beta} \leftarrow \bm{\beta}_i^{(k)}$
       \State $\mathcal{C}_{i-1} \leftarrow \mathcal{C}_{i-1}\cup\{(\alpha,\bm{\beta})\}$
   \EndFor
   \State \textbf{(Forward pass)} for $i=1,\dots,P{-}1$:
          $\hat{\bm{x}}_{i}^{(k+1)} \leftarrow \bm{x}_i^{(k)}$
   \State $\underline{Z}^{(k)} \leftarrow V_1^{(k)}$;\;\;
          $\overline{Z}^{(k)} \leftarrow \sum_{i=1}^{P} F(\mathcal{Y}_i^{(k)})$
   \If{(C1) in \eqref{eq:cv_gap} or (C2) in \eqref{eq:cv_lb} holds}
       \State \textbf{break}
   \EndIf
\EndFor
\State Shut down workers; merge per-window primals into the global trajectory.
\State \Return $\{\mathcal{Y}_i^{(k)}\}_{i=1}^P,\;
                \underline{Z}^{(k)},\;\overline{Z}^{(k)}$.
\end{algorithmic}
\end{algorithm}

\subsection{Discussion}
\label{ssec:discussion}

Three properties distinguish DDDP from overlapping-Schwarz temporal
decomposition and from textbook SDDP. First, the windows are strictly
disjoint, which renders the upper bound $\overline{Z}$ exact rather than
an inner approximation and removes any overlap weight to tune. Second,
the multi-cut, multi-window dual-passing structure generalizes
two-stage Benders to a chain of $P$ stages while remaining fully
parallel within each iteration---unlike SDDP, which sweeps forward and
backward sequentially. Third, the entire temporal coupling is captured
by the boundary dual $\bm{\beta}_i$ of a single fixed initial SOC, so the
cut machinery is agnostic to the OPF model inside the window: any convex
relaxation of the branch-flow or bus-injection formulation---including
the MPISOCP and linear models of Section~\ref{sec:copf}---drops into
$\mathcal{F}_i$ unchanged.
